\section{Introduction}
\label{sec:intro}

A multi-task benchmark scores every model on several sub-datasets and aggregates the scores into one ranking, most often by averaging.
Rankings influence which models are adopted, which papers are published and where funding goes, so developers have reason to act strategically.
One such action is \emph{targeted investment}: folding data related to some sub-datasets into a model's training to raise its scores on those sub-datasets.
Benchmark leakage and contamination are known to inflate scores \citep{zhou2023dont,yang2023rethinking}, and strategic practices distort leaderboards \citep{singh2025leaderboard}, but it is less clear \emph{how much} improvement is needed to change a ranking, and whether the answer depends on how the sub-dataset scores are aggregated.

We study this for a single target model that wants to rank first, with all other models held fixed.
Our question is: what investment makes the target strictly first, and how does it vary with the aggregation rule?
Because ``investment'' has no canonical measure we report three (tasks touched $K$, score points $L_1$, largest single-task raise $L_\infty$) and a convex variant, and because the answer depends on how far a score can plausibly be raised we make the per-task \emph{cap} an explicit modelling assumption and vary it.
Infeasibility, where no plan within the caps reaches first place, is an outcome we report rather than drop.

\paragraph{Contributions.}
\begin{itemize}
\item \textbf{A precise formulation} (Section~\ref{sec:method}) of minimum targeted investment under score-based and rank-based aggregation rules, with exact solvers for nine rules (closed forms or water-filling for score rules, MILPs for rank rules) validated against brute force, and for the mean, median and min against closed-form propositions on every target (Section~\ref{sec:exp-theory}).
\item \textbf{An empirical comparison on public data} (Open LLM Leaderboard, 6 tasks; HELM MMLU, 57 subjects) of how $K$, $L_1$ and feasibility depend on the rule, the target's rank and the cap (Sections~\ref{sec:exp-rules}--\ref{sec:exp-uniform}), with bootstrap intervals over targets.
\item \textbf{Conditions of validity} (Section~\ref{sec:exp-robust}): which conclusions survive raw versus normalised scores, leaderboard slices, task granularity, evaluation noise and leaving one task out, and which do not. Two predictions we made in advance (that the $L_1$ of median and min is of the same order as the mean's or larger, and that ``concentrate on one task'' is the cheapest plan for typical targets) are contradicted or narrowed by the data.
\end{itemize}
The answer we give is therefore conditional: investment is a distribution over rank, rule, cap and cost, and the claims below always carry the cap they were measured under.
We do not model rival reactions, we do not train any model, and the caps are scenario assumptions rather than measurements of what contamination can achieve (Section~\ref{sec:discussion}).
